\documentclass[11pt]{article}
\usepackage[margin=1in]{geometry}
\usepackage{graphicx}
\usepackage{float}
\usepackage{booktabs}
\usepackage{parskip}
\usepackage[hidelinks]{hyperref}

\newcommand{\fig}[1]{../analysis/figures/#1}

% tighten the vertical space around figures
\setlength{\textfloatsep}{10pt plus 2pt minus 2pt}
\setlength{\intextsep}{10pt plus 2pt minus 2pt}

\begin{document}

\begin{flushright}\textbf{Surya Kannan}\end{flushright}

\section*{Summary}
The baseline trains a small language model on Hacker News titles and reaches a
validation loss of \textbf{1.754}. I improved it to \textbf{1.187}, about a
\textbf{32\% reduction}, inside the fixed seven-epoch budget. The key realisation,
shown in the first figure, is that the baseline was not held back by its design
but by how it was trained as the original setup barely trained the model at all.
Everything that followed came from first fixing how the model was trained, a few changes to the training plan and the model, and finally careful
tuning.

\begin{figure}[H]
\centering
\includegraphics[width=\linewidth]{\fig{curves.png}}
\caption{The baseline (red) vs.\ the final model (green). Left is the training
error, right is the validation error. The baseline barely moves, it is close to
where an untrained model starts. The final model keeps improving all the way to
the end of training, and the validation error never turns back up, which means
it is under-trained rather than over-fit.}
\end{figure}

\section{The optimiser}
The original optimiser (SGD) adjusts each part of the model in proportion to
how strong its signal is. In a model like this those signals are very small and
uneven, so the adjustments were tiny and progress stalled. I switched to a more
modern optimiser (AdamW) that automatically rescales each adjustment so the step
size stays sensible no matter how strong or weak the signal is. That single
change is what made the model actually start learning, and it was by far the
biggest improvement of the whole project, from 1.754 down to 1.339.

\section{Training plan and architecture}
With that fixed, a few more well-understood changes each helped a little.

\begin{itemize}
\setlength{\itemsep}{2pt}
\item \textbf{\texttt{warmup\_frac}}: easing the learning speed in gradually before
letting it wind down, rather than starting at full speed, which destabilises early
training (with a cosine wind-down; down to 1.291).
\item \textbf{Initialisation scaling}: a small change to how the model starts out,
so its internal signal stays at a healthy strength as it passes through the layers
(1.242).
\item \textbf{\texttt{lr}}: increasing the learning speed (1.232).
\item \textbf{RoPE}: a more modern way of telling the model where each word sits in
the title. It made better use of the model's existing capacity rather than adding
more, which is why it worked where simply making the model bigger did not (1.205).
\item \textbf{Per-title attention masking}: since the only goal here is to lower
validation loss, one lever I considered was tailoring the model to the specific
kind of data rather than only making generic improvements. This dataset is short,
independent titles, so I stopped attention from crossing from one title into the
next. The effect was neutral, it turned out the model had already learned those
boundaries on its own, a reassuring sign that the idea was right even though it was
redundant.
\end{itemize}

\section{Why bigger didn't help}
Several changes that would normally help did nothing here, like a smaller batch (which
gives more update steps), a deeper model, and a couple of others. The pattern was
consistent, adding more steps or more capacity did not help, and only using the
existing capacity more cleverly did. On top of that, the validation error fell
smoothly to the very end with no sign of the model starting to memorise. The
conclusion is that with only seven epochs the model cannot train enough. It is limited
by the training budget, not by over-fitting.

The clearest evidence came from the warmup experiment. As I increased the warmup,
both the training error and the validation error fell together, in lockstep
(second figure). If the warmup were simply preventing memorisation, the training
error would have risen while the validation error fell. Instead they moved
together, which means the warmup was helping the model train better, not stopping
it from over-fitting.

\begin{figure}[H]
\centering
\includegraphics[width=0.55\linewidth]{\fig{warmup_train_val.png}}
\caption{Training and validation error both fall as the warmup increases, and they
track each other almost exactly. The gap between them does not change, so this is
better training, not reduced over-fitting.}
\end{figure}

\section{Tuning}
I re-tuned the main settings on the final model using a small automated system I
built. The surprise was the warmup. It turned out to be one of the largest levers
in the whole project, and it kept helping far past the usual range. At the higher
learning speed, the very first training steps were doing damage, and giving the
model a long, gentle start fixed it.

That gentle start interacts with the learning speed, since a longer warmup makes a
higher speed safe. So I tested the two together on a small grid (third figure)
instead of one at a time. The best combination, a higher speed paired with a long
warmup, beat anything I would have found by tuning them separately, because the
best speed only works alongside the long warmup. That is the main reason to test
interacting settings together rather than one after another.

Each experiment is just a small configuration file, which is what keeps the whole
process reproducible. The grid above was just a short file.

\noindent\begin{minipage}{\linewidth}
\begin{verbatim}
base:
  weight_decay: 0.1
runs:
  - {name: lr1e-3_w0.4,   lr: 0.001,  warmup_frac: 0.4}
  - {name: lr1.5e-3_w0.6, lr: 0.0015, warmup_frac: 0.6}
  # ...9 learning-speed x warmup combinations in total
\end{verbatim}
\end{minipage}

\begin{figure}[H]
\centering
\includegraphics[width=0.52\linewidth]{\fig{lr_warmup_heatmap.png}}
\caption{Testing learning speed and warmup together (darker is worse, brighter is
better). The best result is a particular \emph{pair} of settings, not simply the
best of each on its own.}
\end{figure}

\section{Engineering}
The single biggest practical lever was iteration speed. The baseline runs on a CPU
at around forty minutes per run, which is far too slow to explore many ideas. I
added a GPU development environment\footnote{\url{https://github.com/MaincodeHQ/mainrun/pull/6}}
that cut a full run down to a couple of minutes, and that speed is what made the
rest possible. With a full sweep taking about fifteen minutes, the cost of checking
an idea is almost nothing. Around that I split the original script into focused
modules, added tests and continuous integration, and built a system that runs
experiments from a configuration file so they stay reproducible.

\begin{figure}[H]
\centering
\includegraphics[width=0.85\linewidth]{\fig{journey.png}}
\caption{The full sequence of experiments. Green points were kept, red were tried
and rejected, and the blue line is the best result so far. The large early drop is
the change of training method; the rest is the training plan, architecture, and tuning.}
\end{figure}

\section{Where I stopped}
This whole approach only works because the runs are cheap. It would not scale to a
real model, where a single training run can cost millions. At that point you cannot
just try everything, so people tune a small stand-in model and carry its settings
over to the big one, predict what will happen from smaller runs, and stop the bad
runs early instead of finishing every one.

I stopped once the gains dropped to tiny fractions. If I kept going, I would not
just search more. The clearest lead is the training plan, the model wanted such a
long warmup that I think the shape of the plan, how the learning speed rises and
falls over training, matters more than the exact numbers I tuned. The architecture
is still open too. RoPE showed that a smarter design can help even when a bigger
model cannot, so I would look at other efficiency-focused changes rather than more
size, and a different optimiser would be worth trying as well.

One honest reflection. Part of why I leaned so heavily on the automated sweeping is
that my grasp of model architecture is still limited. With stronger intuition there
I could have reasoned about which changes would matter most and gone after them
first, rather than building a system to try them and measure. It worked here only
because training is cheap, and the right knowledge would have let me prioritise
where to optimise instead of reaching for hyperparameter tuning by default.

\section*{Final setup and result}
\begin{center}
\begin{tabular}{ll}
\toprule
Training method & AdamW, with weight decay \\
Training plan & higher speed, a long gradual warmup, then a smooth wind-down \\
Architecture & RoPE position handling, per-title attention, tuned initialisation \\
Size & 6 layers, width 512, 16k vocabulary, $\sim$27M parameters \\
Budget & 7 epochs ($\sim$938 steps), batch of 64 \\
\midrule
\textbf{Validation loss} & \textbf{1.186794} (CPU), 1.186144 (GPU) \\
\textbf{vs.\ baseline 1.754} & \textbf{about a 32\% reduction} \\
\bottomrule
\end{tabular}
\end{center}

\end{document}
